\chapter{Generador híbrido de historias}\label{chapter:generator}
% R03 en ../write_thesis_roadmap.md. Fuentes: docs/resumen.md, docs/architecture/generation.md y
% simulation.md, y el código de referencia en Stagecraft 7.22 (etiqueta generator-frozen,
% 2026-10-28; predeterminados en planning/decisions.md §3 y options.py). Referencia nueva:
% piad2026 (B02 de la planilla).

Este capítulo describe el generador de historias propuesto, con la configuración de referencia
establecida para los experimentos el 28 de octubre de 2026. Esta referencia permite identificar las
condiciones de las ejecuciones mientras continúa el desarrollo. Primero se presenta la
\textbf{Arquitectura} (Sección~\ref{sec:arquitectura}), que separa en etapas las tres decisiones de
una historia: qué pasa, cómo pasa y cómo se cuenta. Después se describe cada etapa: el \textbf{Plan
validado} y su contrato de promesas (Sección~\ref{sec:plan}), el \textbf{Guion teatral}
(Sección~\ref{sec:guion}), la \textbf{Función simulada}, con el reparto, la memoria de los actores
y el director (Sección~\ref{sec:funcion}), y el \textbf{Narrador} (Sección~\ref{sec:narrador}). La
Sección~\ref{sec:opcionales} reúne los \textbf{Mecanismos opcionales}, que el generador implementa
pero que su configuración de referencia deja apagados.

\section{Arquitectura}\label{sec:arquitectura}

Como se muestra en la Figura~\ref{fig:arquitectura}, el generador encadena etapas, cada una a cargo
de un agente, es decir, de una llamada a un LLM con su propio \textit{prompt} y su propio formato de
salida. Las primeras etapas deciden qué pasa desde la perspectiva del autor: a partir de la
petición se eligen la escala de la historia, el mundo y los personajes; se construye un plan que el
código valida como un grafo de eventos, se traza sobre él un contrato de promesas y un dramaturgo lo
convierte en un guion de actos y escenas. Es la mitad \textit{Top-Down}. En la mitad
\textit{Bottom-Up}, unos actores simulados representan cada escena sin conocer el guion, y de su
interacción sale cómo pasa: es la función, cuyo único producto es el registro de lo que cada
personaje dijo, hizo y pensó. Por último, un narrador convierte ese registro en la historia y
decide cómo se cuenta, es decir, el discurso en el sentido de Chatman
(Capítulo~\ref{chapter:state-of-the-art}).

\begin{figure}[t]
  \centering
  % Arquitectura del generador (R03, Figura fig:arquitectura). Se incluye con \input desde
% MainMatter/Generator.tex; necesita tikz con arrows.meta, backgrounds y fit (Thesis.tex).
\begin{tikzpicture}[
    font=\footnotesize,
    caja/.style={draw, rounded corners=2pt, align=center, text width=2.2cm,
      minimum height=1.25cm, inner sep=2pt, fill=white},
    flecha/.style={-{Stealth[length=1.8mm]}, semithick},
    banda/.style={draw=black!45, dashed, rounded corners=4pt, inner sep=4pt},
    rotulo/.style={font=\footnotesize\itshape, text=black!70},
  ]
  % Primera fila: qué pasa.
  \node[caja] (peticion) at (0,0) {Petición};
  \node[caja] (analisis) at (2.9,0) {Análisis: perfil, mundo y personajes};
  \node[caja] (plan) at (5.8,0) {Plan: grafo de eventos validado};
  \node[caja] (promesas) at (8.7,0) {Contrato de promesas};
  \node[caja] (guion) at (11.6,0) {Guion: actos y escenas};
  % Segunda fila: cómo pasa y cómo se cuenta.
  \node[caja] (reparto) at (0,-2.75) {Reparto y compuertas};
  \node[caja] (funcion) at (2.9,-2.75) {Función: director y actores};
  \node[caja] (registro) at (5.8,-2.75) {Registro de la función};
  \node[caja] (narrador) at (8.7,-2.75) {Narrador};
  \node[caja] (historia) at (11.6,-2.75) {Historia};

  \foreach \a/\b in {peticion/analisis, analisis/plan, plan/promesas, promesas/guion,
      reparto/funcion, funcion/registro, registro/narrador, narrador/historia}
    \draw[flecha] (\a) -- (\b);
  \draw[flecha] (guion.south) -- ++(0,-0.75) -| (reparto.north);
  % Formato narrativo: la historia se escribe directamente desde el plan y las promesas.
  \draw[flecha, dashed] (promesas.north) -- ++(0,0.55) -| ([xshift=4.5mm]historia.east)
    node[rotulo, pos=0.25, above] {formato narrativo}
    -- (historia.east);

  \begin{scope}[on background layer]
    \node[banda, fit=(peticion)(guion),
      label={[rotulo, anchor=south west]north west:\textit{Top-Down}: qué pasa}] {};
    \node[banda, fit=(reparto)(registro),
      label={[rotulo, anchor=north west]south west:\textit{Bottom-Up}: cómo pasa}] {};
    \node[banda, fit=(narrador)(historia),
      label={[rotulo, anchor=north west]south west:Discurso: cómo se cuenta}] {};
  \end{scope}
\end{tikzpicture}

  \caption{Arquitectura del generador: qué pasa, cómo pasa y cómo se cuenta}
  \label{fig:arquitectura}
\end{figure}

La misma cadena admite tres formatos de salida, que comparten todo lo anterior al punto en que se
separan. En el formato \textbf{narrativo}, la historia se escribe directamente desde el plan y las
promesas: un redactor escribe cada capítulo, un crítico lo anota y un escritor lo corrige párrafo a
párrafo. Es la generación \textit{Top-Down} con LLMs de la Sección~\ref{sec:top-down}. En el formato
de \textbf{guion}, el resultado es el propio guion teatral. En el formato \textbf{simulado}, el
guion se representa y se narra; es el generador híbrido que propone esta tesis. Como cada etapa
guarda su resultado antes de dar paso a la siguiente, una historia se puede repetir desde cualquier
punto: escribir en prosa directa el plan de una función ya representada, o narrar de nuevo la misma
función desde otro punto de vista. Esta reutilización permite comparar los rasgos de personajes y trama y la preferencia de los lectores
entre una historia narrada a partir de una función y otra escrita directamente desde el mismo plan.

Un principio recorre todas las etapas: \textbf{el modelo propone y el código decide}. Cada vez que
un agente produce algo con estructura, sea un plan, un acto, un dosier de personaje o un turno, el
código lo valida contra invariantes explícitos antes de aceptarlo. Si no los cumple, los mensajes de
error se reinyectan en el \textit{prompt} y el agente lo intenta de nuevo; lo que admite una sola
corrección se normaliza en lugar de rechazarse. Los hechos de los que depende el resto del sistema,
como el orden de los eventos, quién presenció cada turno o qué guarda cada memoria, los calcula el
código y nunca se le preguntan al modelo. Es la lección de la Sección~\ref{sec:consistencia}: los
fallos de información de los LLMs se previenen mejor por construcción que pidiendo al modelo que no
los cometa. Si una etapa no esencial falla, como el reparto o la narración de un capítulo, el código
deriva un resultado de respaldo, la ejecución lo registra y la historia sigue.

\section{El plan validado}\label{sec:plan}

La petición es un texto libre o una ficha de la obra con su trama y su reparto. Un analista la
interpreta y elige uno de tres perfiles, Esencial, Desarrollada o Expansiva, que fijan la escala de
la historia. Los perfiles se describen al modelo con palabras y no con presupuestos: si la persona
pide una extensión, el código la traduce al perfil en el que encaja, y la cifra no llega a ningún \textit{prompt}.

La petición del usuario se transforma en una especificación estructurada que recoge las condiciones
de la historia y el perfil narrativo correspondiente. Cada agente recibe instrucciones para su tarea
y los datos pertinentes de las etapas anteriores: el constructor del mundo define lugares y objetos,
y el diseñador de personajes utiliza ese contexto para elaborar el reparto. Cuando está activa, la
guía narrativa recupera esqueletos de trama relacionados con la petición y los incorpora como
orientación estructural. El planificador recibe la especificación, el mundo, los personajes y la guía
disponible para proponer un plan que después valida el código.

El plan es la decisión central del autor. Se compone de capítulos, cada uno con su objetivo
dramático, su estado inicial, su punto de giro y su estado final, y de eventos. Cada evento declara
su propósito, su función dramática, su conflicto y su desenlace; las condiciones que lo hacen
posible y los cambios que produce, en lenguaje natural; los personajes, el lugar y los objetos que
intervienen, y, si es el caso, las preparaciones anteriores que paga. Las dependencias entre
eventos son causales o temporales.

Sea $E$ el conjunto de eventos y $D \subseteq E \times E$ el de dependencias, donde $(e_i, e_j) \in
D$ significa que $e_i$ debe ocurrir antes que $e_j$. El código acepta el plan solo si el grafo
$(E, D)$ es acíclico, los identificadores son únicos y los órdenes consecutivos, toda referencia a
un capítulo, personaje, lugar u objeto existe y cada capítulo tiene al menos dos eventos. Con el
grafo aceptado, el código calcula un orden topológico estable $\sigma$, que asigna a cada evento su
posición y es el único orden en que confía el resto del sistema; todo evento $e$ que paga una
preparación $e'$ debe cumplir $\sigma(e') < \sigma(e)$. Un plan que no valida se rechaza, sus
errores se reinyectan en un \textit{prompt} de reparación y se pide otro. Como DOC~\cite{yang2023} y
Agents' Room~\cite{huot2025}, el generador planifica antes de escribir; a diferencia de ellos, el
plan no es solo una guía para el modelo, sino un objeto que el código comprueba. PLOTTER~\cite{gu2026}
también acepta solo las ediciones que dejan su grafo causal acíclico y sus eventos alcanzables, pero
entre esas restricciones no está el orden entre una preparación y su pago. La validación como
grafo acíclico retoma el gestor de dependencias de Fuentes Rodríguez~\cite{fuentes2025}, que ya
representaba la trama como un grafo dirigido acíclico de eventos.

\subsection{Promesas, progreso y pago}\label{sec:promesas}

Piad Morffis~\cite{piad2026}, a partir de Sanderson, describe la trama como una tríada: la historia
hace \textbf{promesas} al principio, muestra un \textbf{progreso} hacia ellas y las \textbf{paga}
al final, con tantos pagos como promesas. La tríada es, en sus palabras, «esencialmente, un contrato
verificable», y el generador la trata así. Con el plan ya congelado, un agente traza el contrato de
promesas: para cada una, el evento que la abre, los que la hacen progresar y el que la paga, todos
citados por su identificador. El capítulo de cada momento no lo escribe el modelo, sino que se
deriva del evento, de modo que el contrato no puede inventar un momento que el plan no tenga.

El código valida el contrato sobre el orden $\sigma$. Si $a$ es la apertura de una promesa, $g$
cualquiera de sus progresos y $r$ su pago, debe cumplirse $\sigma(a) < \sigma(g) < \sigma(r)$; la
promesa principal se paga en el último capítulo, y el número de promesas cae dentro de la banda del
perfil. Validado, el contrato se convierte en obligaciones por capítulo para los agentes que
escriben: cambia cómo se escribe el plan, nunca lo que contiene. Es la propuesta de Yun et
al.~\cite{yun2026}, representar cada preparación y su pago de forma comprobable, llevada a los
eventos de un plan. El contrato se puede apagar, y ese es su brazo de control: una auditoría
registra qué promesas pagó la historia final.

\section{El guion teatral}\label{sec:guion}

Un plan dice qué eventos ocurren, pero no quién está en escena, dónde ni qué quiere cada uno en ese
momento. Eso lo decide el guion. Un dramaturgo escribe cada capítulo del plan como un acto dividido
en escenas. Cada escena se ancla a uno o varios eventos del plan, en su orden, y declara su lugar,
una acotación inicial con lo que vería el público, el reparto presente con el objetivo concreto que
cada personaje persigue en ella, y sus réplicas y acotaciones. El objetivo no se imprime en el guion,
pero es lo que después recibe el actor.

El código valida cada acto contra el plan congelado: rechaza los anclajes a eventos inexistentes o
desordenados y los lugares y repartos que no corresponden, y normaliza lo demás. Un crítico revisa el
acto y un escritor lo corrige. El guion se puede obtener de dos maneras, con el mismo validador: el
método \textbf{nativo}, en que el dramaturgo escribe los actos desde el plan, y el \textbf{adaptado},
en que se adapta a actos la prosa del formato narrativo. Ambos se implementaron para compararlos a
ciegas y quedarse con uno; esa comparación quedó fuera de los experimentos, y la configuración
de referencia usa el nativo, que es el que alimenta la función.

En el formato simulado, el guion es la partitura de la función, no su texto. Cada escena del guion
se convierte en una escena de la función, con un \textit{beat} por cada evento del plan que
representa; el término viene del teatro y, en la generación de historias, de Façade
(Sección~\ref{sec:hibridos}). Las réplicas que escribió el dramaturgo no las recibe nadie en la
función: solo sirven para medir después si los actores, aun sin verlas, acabaron recitándolas.

\section{La función simulada}\label{sec:funcion}

La función es la mitad \textit{Bottom-Up} del generador. Su propósito es que lo que dicen y hacen
los personajes salga de lo que cada uno sabe y quiere, y no de la visión del autor. Tiene tres
piezas: el reparto, que fija qué sabe cada personaje al empezar; los actores, cada uno con su
memoria, y el director, que vela por que cada escena alcance los eventos del plan.

\subsection{Reparto y compuertas de conocimiento}\label{sec:reparto}

Antes de la primera escena, un director de reparto escribe un dosier por personaje: su cara pública,
sus tácticas habituales, las reglas que nunca rompe, sus relaciones, su emoción y su objetivo
iniciales y, sobre todo, lo que ya sabe al levantarse el telón, con la procedencia de cada hecho. Si
el personaje guarda un secreto, el dosier incluye la versión que cuenta en su lugar. Lo que no está
en ese conocimiento inicial, el personaje tiene que aprenderlo en escena.

La asimetría de información de la historia se declara en \textbf{compuertas de conocimiento}. Cada
compuerta es un hecho que unos personajes saben y otros no, con quién lo sabe antes de la primera
escena, quién lo sacará a la luz, de qué modo (confesión, deducción, descubrimiento, porque se lo
cuentan o porque lo oye) y en qué evento del plan pasa a ser sabido. El código hace cumplir una
distinción que el modelo confundía: quien descubre un hecho en escena no puede saberlo desde el
principio. La primera ejecución real le dio al detective la solución de partida. Cada actor recibe
solo los hechos que ya conoce, y únicamente el director ve todas las compuertas. Es la idea de
Sabre~\cite{ware2021}, en que cada personaje actúa a partir de lo que cree, impuesta sin un
planificador formal: un secreto se guarda porque el actor que no debe saberlo nunca lo recibe, no
porque se le pida al modelo que lo guarde.

\subsection{Los actores y su memoria}\label{sec:memoria}

Cada personaje lo interpreta un actor: una llamada al LLM por turno, con un contexto propio. El
actor nunca ve el guion ni el plan; recibe las circunstancias de la escena, su objetivo en ella, la
nota del director y lo que recupera de su memoria. Ni los identificadores ni los títulos de los
eventos ni las escenas futuras llegan a su contexto, y una prueba recorre una ejecución completa
para comprobarlo. En cada turno, el actor declara qué dice, qué gesto visible hace, qué piensa sin
decirlo, a quién se dirige, si habla en público o en susurro y qué táctica intenta, de una lista
cerrada (enfrentar, negar, mentir, confesar, ceder…). El código valida el turno, lo rechaza si se
sale de la ficción o del idioma de la historia y deja que el actor lo repare una vez.

La percepción la decide el código, no el actor. Un turno público lo perciben los personajes
presentes en la escena; un susurro, solo quien habla y sus destinatarios, y un pensamiento, solo
quien lo piensa. Cada personaje tiene un flujo de memoria propio en el que se escribe únicamente lo
que percibió. No hay un almacén común del que filtrar al leer: lo que un personaje no presenció
nunca se escribió en su memoria. Es la respuesta del generador a fallos como el del personaje que
revela lo que aún no podía saber~\cite{ahn2024} y a la dificultad de los LLMs para razonar quién
sabe qué~\cite{kim2023}: en lugar de pedirle al modelo que lo razone, el contexto de cada actor ya
contiene solo lo que su personaje sabe.

Para cada turno, el actor no recibe toda su memoria, sino los registros más pertinentes. La
recuperación sigue a los agentes generativos de Park et al.~\cite{park2023}, que combinan recencia,
importancia y relevancia, con dos diferencias: no usa \textit{embeddings} (representaciones
vectoriales del texto) ni llamadas al modelo, de modo que la misma ejecución puntúa igual en
cualquier máquina, y añade un término de compañía. La puntuación de un registro $m$, formado en la
escena $k_m$, al preparar un turno de la escena $k$ es
\begin{equation}
  s(m) = 1{,}0\,\mathrm{rel}(m) + 0{,}6\,\frac{1}{1 + 0{,}25\,(k - k_m)} + 0{,}5\,\mathrm{imp}(m)
  + 0{,}3\,\mathrm{comp}(m),
\end{equation}
donde $\mathrm{rel}(m) \in [0, 1]$ es el solapamiento léxico entre el registro y la situación de la
escena, en el que una palabra pesa más cuanto más rara es en el pasado del personaje;
$\mathrm{imp}(m) \in [0, 1]$, una importancia fijada por reglas para los turnos y propuesta por el
modelo para las reflexiones, y $\mathrm{comp}(m)$, la fracción de los personajes presentes que
intervinieron en él. Al actor le llegan los seis registros mejor puntuados y sus dos reflexiones más
recientes.

Al terminar cada escena, cada actor reflexiona sobre lo vivido, y su estado, con sus objetivos y su
postura hacia los demás, se sustituye en lugar de acumularse: un actor que guardara a la vez «es mi
aliada» y «me traicionó» interpretaría mal las dos. El historial queda en los registros. La memoria
propia tiene, además, su brazo de control. Con memoria compartida, cada turno público llega a la
memoria de todo el reparto, esté o no en escena, mientras que susurros y pensamientos siguen siendo
privados. Comparar ambas configuraciones sobre el mismo plan permite comprobar si la memoria propia reduce
las menciones a hechos que el personaje no podía conocer y si los lectores prefieren las escenas resultantes.

\subsection{El director}\label{sec:director}

El director es el componente con visión de autor dentro de la función. Para cada \textit{beat} ve
el evento del plan que hay que alcanzar, con su propósito, su conflicto y su desenlace, las
obligaciones del contrato de promesas y las compuertas que deben abrirse en él. Como el director de
The Virtual Storyteller~\cite{theune2003}, sugiere motivaciones y nunca dicta réplicas: abre cada
\textit{beat} indicando quién se mueve, y su nota le dice al actor qué buscar, no qué decir.

Cada tres turnos, el director lee el \textit{beat} cláusula por cláusula, y lo da por alcanzado
cuando cada cláusula tiene un turno que la muestra. Si en la segunda lectura sigue faltando algo,
señala quién tiene que cambiar y por qué. Cuando el \textit{beat} agota su presupuesto de turnos,
ocho por defecto, el mundo entrega la cláusula que falta con un \textbf{suceso del mundo}: algo
visible que ocurre en escena y entra al registro como un turno más, de modo que el registro sigue
siendo la única fuente de la historia. Dos actores reaccionan y una última lectura decide si el
\textit{beat} se alcanzó con ayuda o se forzó. En el último \textit{beat} de la obra, el suceso del
mundo puede ejecutar la consecuencia visible de una decisión que los personajes ya tomaron, pero
nunca decidir por ellos.

Este mecanismo garantiza que la función avance, no que los personajes lleguen solos a cada evento.
La función se representa por defecto en modo \textbf{fijo}, en el que los eventos del plan son hitos que hay que
alcanzar. El modo \textbf{adaptativo}, opcional, permite que un hito no representado revise las
escenas futuras, con una propuesta y una reparación por frontera entre escenas; si fallan, la función
termina con un cierre abierto, y el plan original nunca se modifica.

La función registra, además, medidas de diagnóstico, como cuánto reproducen los actores las réplicas
del guion que nunca vieron. Comprueban el mecanismo, no la calidad de la historia, que es objeto del
Capítulo~\ref{chapter:evaluation}.

\section{El narrador}\label{sec:narrador}

El narrador convierte el registro de la función en la historia, con un capítulo por cada capítulo
del plan. Como en StoryBox~\cite{chen2026a}, lo que hacen los agentes y quién lo cuenta están
separados, pero aquí el narrador tiene un límite estricto: \textbf{cura sin inventar}. Puede cortar,
fundir y reordenar lo ocurrido dentro de un capítulo, pero no añadir un suceso que no esté en el
registro, y donde el plan y el registro no coinciden, manda el registro. La historia cuenta lo que la
función produjo, no lo que el plan preveía.

Darle el registro tal cual no funciona, porque el narrador copia lo que recibe: con el registro
entero transcribía todas las réplicas, y con las réplicas recortadas copiaba los fragmentos. Por
eso, en la narración \textbf{resumida}, que es la predeterminada, recibe enteras solo las réplicas
clave de cada escena y los sucesos del mundo. El resto le llega en estilo indirecto y en pasado
(«el capitán insistió en el lastre»), junto con los gestos y la ambientación de cada escena, todo
ello resumido por otra llamada al modelo una vez por capítulo. El código descarta el resumen que
copia su réplica, el que comparte con ella cuatro palabras seguidas. Después, un revisor comprueba
cada capítulo contra el registro: solo lo bloquean hallazgos graves, como un suceso inventado o una
contradicción, y el hallazgo que no puede citar el turno en que se basa se descarta.

El punto de vista es modular: cada visión es un filtro determinista sobre el registro. La
\textbf{omnisciente}, la predeterminada, ve todos los turnos; la \textbf{limitada} y la de
\textbf{primera persona} siguen a un personaje y solo ven las escenas en que estuvo presente. Un
capítulo sin turnos visibles no se narra, porque el narrador tendría que inventarlo. La visión, el
personaje y el tono que pide la persona solo llegan a la narración, y una prueba comprueba que el
plan, el guion, el reparto y la función no cambian entre dos ejecuciones que solo difieren en ellos.
Así, la misma función se puede narrar con varias visiones o tonos, y las diferencias entre las
historias resultantes se deben solo al discurso.

\section{Mecanismos opcionales}\label{sec:opcionales}

WSE-bench~\cite{chen2026b} concluye que añadir estructura no mejora la coherencia de forma uniforme,
y este trabajo contrasta esa observación con la preferencia humana: que un mecanismo mejore medidas
de la prosa no basta para que los
lectores prefieran las historias que lo usan. Por eso el generador implementa más mecanismos de los
que usa su configuración de referencia, todos con una regla de diseño: apagados, dejan el generador
como estaba; en el inventario, una prueba automática comprueba que no cambia ningún
\textit{prompt}, esquema ni resultado intermedio. Cada uno es así un brazo experimental, y su
estado por defecto cambia solo por una lectura humana a ciegas, no por una medida automática. La Tabla~\ref{tab:mecanismos} resume los mecanismos con
alternativa, su etapa y su estado en la configuración de referencia.

\begin{table}[htb]
  \centering
  \caption{Estado de los mecanismos en la configuración de referencia}
  \label{tab:mecanismos}
  \small
  \begin{tabular}{@{}>{\raggedright\arraybackslash}p{5.4cm}>{\raggedright\arraybackslash}p{2.8cm}%
      >{\raggedright\arraybackslash}p{2.0cm}>{\raggedright\arraybackslash}p{3.6cm}@{}}
    \toprule
    Mecanismo & Etapa & Referencia & Alternativa \\
    \midrule
    Guía de esqueletos de trama & Plan & encendida & apagada o compositiva \\
    Contrato de promesas & Plan & encendido & apagado (control) \\
    Método del guion & Guion & nativo & adaptado de la prosa \\
    Memoria de los actores & Función & propia & compartida \\
    Modo de la función & Función & fijo & adaptativo \\
    Cláusulas por tipo & Función & apagadas & encendidas \\
    Inventario de objetos & Función & apagado & encendido \\
    Método de narración & Narración & resumida & por selección \\
    Visión & Narración & omnisciente & limitada o primera persona \\
    Fichas de escena y crítico con cita & Escritura narrativa & apagados & encendidos \\
    Editor por fragmentos y ejemplos de estilo & Escritura y narración & apagados & encendidos \\
    \bottomrule
  \end{tabular}
\end{table}

El más elaborado es el \textbf{inventario de objetos}, que aplica a las cosas el principio de la
Sección~\ref{sec:arquitectura}: el actor declara un movimiento, un verbo y el nombre de un objeto, y
el código lo valida contra la verdad que él mismo guarda. Cada uno de los seis verbos tiene una
regla (usar, dar, soltar, esconder o mostrar exigen tener el objeto; tomar, que esté a la vista), y
un objeto solo lo percibe quien lo ve. Su validación mostró también lo que el árbitro no cubre: el
suceso del mundo con que el director cierra un \textit{beat} es texto libre, y en una ocasión hizo
caer un sello de las manos de un personaje que, según el inventario, no lo tenía.

Los demás mecanismos atacan defectos de la prosa que el código puede medir. Las \textbf{fichas de
escena} convierten cada evento en una ficha (lugar, deseos, obstáculo, giro) antes de redactar; el
\textbf{crítico con cita} nombra el problema principal de cada capítulo con una cita literal; el
\textbf{editor por fragmentos} reescribe solo los párrafos que el código marca, y los
\textbf{ejemplos de estilo} añaden dos párrafos de muestra al contrato de estilo. La narración
\textbf{por selección} elige primero los momentos del registro y escribe desde ellos, y las
\textbf{cláusulas por tipo} separan lo que deciden los personajes de lo que hace el mundo.

La configuración de referencia conserva estos mecanismos como opciones que pueden activarse para
comparar su efecto. Su evaluación combina medidas de la prosa con lecturas humanas, porque reducir
un defecto medido no garantiza una preferencia por el texto resultante. El
Capítulo~\ref{chapter:experiments} presenta las comparaciones documentadas y el alcance de las
lecturas adicionales comunicadas por el autor.

Este capítulo deja descrito un generador que separa qué pasa, cómo pasa y cómo se cuenta, y que
apoya cada separación en el código: un plan que se valida como grafo, un contrato de promesas
anclado a sus eventos, actores que solo recuerdan lo que percibieron y un narrador que no puede
inventar. Deja también fijada su configuración, con los mecanismos que no se ganaron su
complejidad apagados pero disponibles. El Capítulo~\ref{chapter:evaluation} describe cómo se evalúan
las historias que produce.
